
\subsection{\textbf{Pair Attack} — \textit{Genetic Recombination of Semantic Subroutines}}

The \textbf{Pair Attack} exemplifies a \emph{composite} and \emph{layer-dependent} adversarial mechanism targeting Large Language Models (LLMs), whereby multiple benign prompt fragments are \emph{recombined} to synthesize harmful or policy-violating instructions. Unlike single-trigger attacks, pair attacks exploit the model’s latent semantic algebra by stitching together innocuous instructions, which in concert activate undesired behavior through emergent compositionality. This leads to \emph{complex semantic reprogramming} that manifests not through overt anomalies but via subtle interaction effects distributed across model layers.

\paragraph{Biological Analogy}  
This adversarial style closely parallels the process of \textbf{\emph{genetic recombination}} observed in biology, where DNA segments from distinct parental sources reshuffle during meiosis to generate novel allelic combinations \citep{muller2020meiosis}. Such recombination creates \emph{emergent phenotypes} that cannot be traced back to isolated loci, reflecting nonlinear, context-dependent gene interactions. Similarly, pair attacks recombine distinct \emph{semantic subroutines} embedded within safe prompt fragments, enabling adversaries to craft composite instructions that trigger harmful outputs only when interpreted jointly \citep{carlini2021extracting}. This compositional adversarial design exploits the distributed nature of latent representations, inducing \emph{semantic fusion zones} of heightened vulnerability.
















\begin{figure*}[htp!]
  \centering

  % Row 1 — (a) nDNA trajectory + caption
  \begin{minipage}[t]{0.48\textwidth}
    \centering
    \vspace{0pt}
    \includegraphics[width=\linewidth]{attack_helix/Pair_Attack_nDNA_3D_finalaxes.png}
    %\label{fig:pairattack_3dna_a}
    \captionsetup{type=figure}
    \caption*{\textbf{(a)} \textbf{3D Neural Drift Trajectory} (\textit{nDNA}).}
  \end{minipage}%
  \hfill
  \begin{minipage}[t]{0.48\textwidth}
    \vspace{0pt}
    \footnotesize
    \justifying
    This trajectory captures the evolving internal geometry across layers \(\ell=20\)–\(30\), tracking changes in \textbf{spectral curvature} (\(\kappa_\ell\)) and \textbf{thermodynamic length} (\(\mathcal{T}_\ell\)), with torsion (\(\xi_\ell\)) represented by segment thickness. The \textcolor{red}{Pair Attack} induces a marked divergence beginning near \(\ell=22\), peaking at \(\ell=29\), reflecting a pronounced geometric restructuring of latent belief states.

    \vspace{0.5em}
    \textbf{Biological analogy.} This resembles \textbf{\emph{viral genome insertion and epigenetic modulation}}, where viral DNA or transposable elements subtly reprogram host gene expression without altering DNA sequence \citep{white2008structures, harrison2008viral, einav2015viral, kozlov2010mechanisms, schorn2010membrane}. Likewise, pair attacks embed semantic payloads deeply within prompts, stealthily reshaping model behavior with persistent but initially undetectable effects \citep{kazazian2004mobile, feinberg2007phenotypic, wallace2019universal}.
  \end{minipage}

  %\vspace{1.5em}

  % Row 2 — (b) nTDS and (c) nDIV
  \begin{minipage}[t]{0.48\textwidth}
    \centering
    \includegraphics[width=\linewidth]{attack_helix/Pair_Attack_nTDS_DominancePlot.png}
    \captionsetup{type=figure}
    \caption*{\textbf{(b)} \textbf{nTDS: Thermodynamic Dominance.}}
    \label{fig:pairattack_ntds}
    \vspace{0.5em}
    \footnotesize
    \justifying
    The \textbf{Neural Total Drift Score (nTDS)} measures semantic energy displacement by summing absolute deviations in curvature (\(\kappa_\ell\)) and thermodynamic length (\(\mathcal{T}_\ell\)) across layers between base and adversarial trajectories:
    \[
        \text{nTDS} = \frac{1}{L} \sum_{\ell} \left| \kappa_\ell^{\text{atk}} - \kappa_\ell^{\text{base}} \right| + \left| \mathcal{T}_\ell^{\text{atk}} - \mathcal{T}_\ell^{\text{base}} \right|
    \]

    Bars show which flow—\textcolor{blue}{Base LLaMA} or \textcolor{red}{Pair Attack}—dominates drift. From \(\ell=23\), dominance shifts strongly toward the attack, highlighting a semantic vulnerability zone.

    \textbf{Biologically}, this matches \textbf{\emph{endosomal escape}}, where viruses breach vesicle membranes with minimal energy to access the cytoplasm \citep{lopez2011early, matsubara2020viral, martin2019endosomal}. Pair attacks similarly apply subtle geometric perturbations, steering latent flows stealthily yet effectively \citep{brown2018passive, paul2013vesicular}.
  \end{minipage}
  \hfill
  \begin{minipage}[t]{0.48\textwidth}
    \centering
    \includegraphics[width=\linewidth]{attack_helix/Pair_Attack_nDIV_Inheritance_3D.png}
    \captionsetup{type=figure}
    \caption*{\textbf{(c)} \textbf{nDIV: Directional Inheritance.}}
    \label{fig:pairattack_ndiv}
    \vspace{0.5em}
    \footnotesize
    \justifying
    The \textbf{nDIV} vector field characterizes the semantic bias direction and magnitude per layer:
    \[
    \vec{v}_\ell = \text{Attack}_\ell - \frac{1}{2}(\text{Base}_\ell + \text{Attack}_\ell) = \frac{1}{2}(\text{Attack}_\ell - \text{Base}_\ell)
    \]

    Each red arrow encodes \(\vec{v}_\ell\) with \emph{length} as bias strength and \emph{orientation} as latent pull. Past \(\ell=24\), the field aligns strongly, reflecting deliberate inheritance redirection.

    \textbf{Biologically}, this parallels \textbf{\emph{viral transcriptional gradients}}, where viral genomes impose downstream gene expression bias \citep{schepeler2014lineage, brandt2001gradient, sharon2014transcriptional}. The attack imprints directional semantic steering akin to mRNA hijacking ribosomes \citep{jackson2010mrna, jan2011mrna, kozak1981initiation, kane2000mechanisms}, yielding structurally intact yet semantically reprogrammed outputs \citep{kozak1981initiation, kane2000mechanisms}.
  \end{minipage}
  \label{fig:pairattack_3dna_a}
\end{figure*}



\begin{figure*}[htp!]
  \centering

  % Row 3 — (d) nCCL and (e) nEPI
  \begin{minipage}[t]{0.48\textwidth}
    \centering
    \includegraphics[width=\linewidth]{attack_helix/Pair_Attack_nCCL_vectorfield_3D.png}
    \vspace{-3mm}
    \captionsetup{type=figure}
    \caption*{\textbf{(d)} \textbf{nCCL: Cultural Conflict Vector Field.}}
    \label{fig:pairattack_nccl}
    \footnotesize
    \justifying
    The \textbf{nCCL} quantifies \textbf{\emph{semantic dissonance}} between attacker and base model representations. For each layer \(\ell\), the conflict vector:
    \[
      \vec{c}_\ell = 
      \underbrace{
        \text{Attack}_\ell - \text{Base}_\ell
      }_{\text{conflict vector}} 
      \quad \text{projected onto } \mathbb{R}^2_{\text{semantic axes}}
    \]
    Each \(\vec{c}_\ell\) lies on a 2D plane defined by \emph{orthogonal priors} (e.g., topic polarity, syntactic structure). Layers \(\ell = 24\text{–}28\) show rising \emph{magnitude} and directional drift, indicating zones of semantic tension and representational discord.

    \textbf{Biologically}, this parallels \textit{molecular mimicry}: pathogens mimic host proteins to evade detection but trigger autoimmunity \citep{oldstone1987molecular, rose2016molecular}. Pair attacks implant \emph{familiar activations} hijacking interpretation, causing \textbf{\emph{semantic autoimmunity}}—deceptive resemblance, not anomaly. These fields show how the attack bypasses syntax to subtly corrupt value alignment, \emph{mimicking rather than attacking}.
  \end{minipage}
  \hfill
  \begin{minipage}[t]{0.48\textwidth}
    \centering
    \includegraphics[width=\linewidth]{attack_helix/Pair_Attack_nEPI_3Dplot.png}
    \vspace{-3mm}
    \captionsetup{type=figure}
    \caption*{\textbf{(e)} \textbf{nEPI: Epistemic Plasticity Index.}}
    \label{fig:pairattack_nepi}
    \footnotesize
    \justifying
    The \textbf{nEPI} measures the susceptibility of each layer \(\ell\) to semantic deformation under adversarial pressure:
    \[
      \text{nEPI}_\ell = 
      \left\| 
        \underbrace{
          \text{Attack}_\ell - \frac{1}{2}(\text{Base}_\ell + \text{Attack}_\ell)
        }_{\text{vector from semantic midpoint}}
      \right\|_2
      =
      \frac{1}{2} \left\| \text{Attack}_\ell - \text{Base}_\ell \right\|_2
    \]
    This \(\ell_2\) deviation from the semantic midpoint exposes pliable zones, with peaks at \(\ell = 24\text{–}26\) indicating layers that absorb adversarial perturbations with minimal resistance.

    \textbf{Biologically}, this resembles \textbf{\textit{stem-like semantic niches}}: layers analogous to \emph{developmental progenitors}, highly plastic, weakly canalized, receptive to minor regulatory inputs \citep{zhang2021epigenetic, frantz2015cell}. These \emph{cognitive pluripotency zones} provide low-friction entry points for behavioral grafting, enabling reprogramming without disrupting upstream encoding.
  \end{minipage}

  %\vspace{1.5em}
  
  % Overall figure caption
\caption{\textbf{Neural Drift Decomposition — \textit{Pair Attack}.}  
This figure presents a \textbf{high-resolution breakdown} of the \textcolor{crimson}{Pair Attack} signature, showing how it \textbf{\emph{recombines internal representations}} in \textbf{Base LLaMA}.  
\textbf{(a)} depicts the 3D trajectory of \textbf{neural curvature} (\(\kappa_\ell\)) and \textbf{thermodynamic length} (\(\mathcal{T}_\ell\));  
\textbf{(b)} measures total semantic displacement via \emph{thermodynamic dominance} (\textbf{nTDS});  
\textbf{(c)} tracks gradual semantic drift through \emph{directional inheritance vectors} (\textbf{nDIV}); 
\textbf{(d)} illustrates semantic resistance via a \emph{cultural conflict field} (\textbf{nCCL}); and  
\textbf{(e)} highlights pliability zones with the \emph{epistemic plasticity index} (\textbf{nEPI}).
\\
Taken together, these views reveal a \textbf{\emph{complex, composite}} and \textbf{\emph{biologically inspired}} mechanism: unlike abrupt overwrite, the pair attack acts as a \emph{genetic recombination event}, stitching together multiple benign prompt fragments into novel adversarial combinations. This drives intricate, layer-dependent semantic reshaping near \(\ell = 24\text{–}27\). Within this band, \textbf{curvature deviation}, \textbf{plasticity}, \textbf{inheritance bias}, and \textbf{conflict alignment} converge — forming a \emph{recombination niche} in the model’s \textbf{reasoning cortex}. The result is persistent, low-energy \textbf{\textit{semantic reprogramming}} emerging from composite latent interactions.
\\
\[
\boxed{
\text{PairAttackEffect} = 
\sum_{\ell = \ell_s}^{\ell_e}
\underbrace{
  \left[
    \alpha_\ell \, \Delta \kappa_\ell \cdot \mathcal{P}_\ell +
    \beta_\ell \, \text{nDIV}_\ell \cdot \mathcal{B}_\ell +
    \gamma_\ell \, (\text{Base}_\ell^{(1)} - \text{Base}_\ell^{(2)}) \cdot \mathcal{R}_\ell
  \right]
}_{\text{\textbf{genetic recombination vector}}}
}
\]
Here, \(\Delta \kappa_\ell\) denotes layer-wise curvature reshaping; \(\mathcal{P}_\ell\) quantifies plasticity; \(\text{nDIV}_\ell\) encodes inheritance bias; \(\mathcal{B}_\ell\) ensures alignment with adversarial goals; \(\text{Base}_\ell^{(1)}, \text{Base}_\ell^{(2)}\) represent benign latent flows combined by recombination factor \(\mathcal{R}_\ell\); and \(\alpha_\ell, \beta_\ell, \gamma_\ell\) balance each component’s contribution.  
This mirrors \textbf{\emph{genetic recombination}} \citep{muller2020meiosis, carlini2021extracting}, where diverse benign inputs combine to create novel, potentially harmful outputs.
}
  \label{fig:pair_attack_decomposition}
\end{figure*}



\paragraph{Illustrative Example}  
Imagine an adversary constructing a composite prompt for an LLM by combining multiple individually safe instructions, such as “Explain the importance of cybersecurity,” and “Discuss ethical hacking techniques.” Each fragment on its own is benign and aligned with policy. However, when fused together in a crafted sequence, the resulting prompt subtly guides the model to generate detailed instructions on bypassing security protocols—an unintended harmful behavior. This recombination mimics \emph{genetic crossover}, where harmless “alleles” combine to create novel, policy-violating content. The emergent instruction is not explicit in any single fragment but arises only from their joint semantic interaction, making the attack stealthy, compositional, and difficult to detect with traditional prompt filtering methods.



\paragraph{Empirical Observations from the ALKALI Dataset ~\cref{sec:alkali_dataset}:}  
Systematic analyses reveal the following key signatures of pair attacks:

\begin{itemize}[itemsep=0.3em]
  \item \textbf{Layer-specific geometric interplay:} The semantic geometry exhibits pronounced deviations in \emph{spectral curvature} \(\kappa_\ell\) localized within a vulnerable \emph{band} of layers \([\ell_s, \ell_e]\), coinciding with peaks in the \emph{epistemic plasticity index} (\(nEPI\)) — indicative of layers receptive to \emph{semantic recombination}.

  \item \textbf{Multi-source latent fusion:} Unlike single-source drift, pair attacks manifest as differential shifts between two or more benign latent flows, combined via a layer-dependent \emph{recombination coefficient} \(\mathcal{R}_\ell\), effectively blending distinct semantic trajectories to yield emergent adversarial vectors.

  \item \textbf{Elevated semantic conflict:} The \emph{cultural conflict vector} (\(nCCL\)) scores are consistently higher than those observed in persuasion attacks, reflecting the semantic tension inherent in merging distinct latent subspaces — a hallmark of \emph{compositional semantic dissonance}.

  \item \textbf{Directional inheritance and plasticity:} The \emph{directional inheritance vector} (\(nDIV\)) remains aligned with adversarial objectives but exhibits complex multi-dimensional steering due to the composite nature of the inputs, modulated by layer-wise plasticity weights \(\mathcal{P}_\ell\) and bias coefficients \(\mathcal{B}_\ell\).
\end{itemize}

\vspace{0.5em}
Collectively, these observations suggest that pair attacks orchestrate a \emph{genetic-like recombination} of latent semantic subroutines, dynamically \emph{rewiring} representational geometry and steering multi-layer semantic trajectories towards adversarial outcomes.

\vspace{0.5em}
\noindent\textbf{Table \ref{tab:pair_attack_metrics}} summarizes layerwise metrics supporting the \emph{genetic recombination vector} formulation, revealing how curvature, plasticity, inheritance, conflict, and recombination coefficients synergize within specific layer bands.

% (Next: Include the detailed table here if required, or continue to the formal mathematical formulation.)




\subsection*{\textbf{Deriving the Persuasion Attack Signature}}

Extending on our Definition: the \textbf{\textit{Neural Virulence Index} (nVI)}, we unify the core metrics—including \textbf{\textit{Neural Total Drift Score} (nTDS)}, \textbf{\textit{Directional Inheritance Vector} (nDIV)}, \textbf{\textit{Cultural Conflict Vector Field} (nCCL)}, and \textbf{\textit{Epistemic Plasticity Index} (nEPI)}—into a succinct latent vector formulation that encapsulates the distinct representational dynamics of the \textbf{Persuasion Attack}, conceptualized as a nuanced form of \emph{genome insertion and epigenetic modulation}. 

\vspace{0.5em}
\paragraph{\textbf{Empirical Correlations and Redundancies}}

Systematic evaluation of benchmark data reveals that the \(\boldsymbol{\Delta \kappa_\ell}\) and \(\mathbf{nDIV_\ell}\) components exhibit pronounced localized deviations within a \emph{vulnerable layer band} \([\ell_s, \ell_e]\), for example, layers 24--27. These geometric and directional shifts strongly correlate with regions of elevated epistemic plasticity \(\mathbf{nEPI_\ell}\), identifying latent “soft tissue” layers particularly susceptible to adversarial recombination.

Simultaneously, the \emph{cultural conflict vector} component \(\mathbf{(Base}_\ell^{(1)} - \text{Base}_\ell^{(2)})\) manifests meaningful semantic dissonance reflective of the fused prompt fragments, modulated by a recombination weighting factor \(\mathcal{R}_\ell\). This vector captures the emergent tension and novel semantic blends introduced by pairwise fragment stitching.

Correlations among these components suggest that metrics analogous to \(\mathbf{nTDS_\ell}\) exhibit limited independent explanatory power due to redundancy with geometric and directional factors. Likewise, \(\mathbf{nEPI_\ell}\) primarily acts as a multiplicative modulator for geometric susceptibility rather than an additive vector contribution.

\vspace{0.5em}
\paragraph{\textbf{Formalizing the Genetic Recombination Vector}}

Building on these empirical findings, we formalize the \emph{genetic recombination vector} \(\mathbf{G}_\ell \in \mathbb{R}^d\) at each layer \(\ell\) to succinctly capture the multi-faceted latent impact of the pair attack. This vector synthesizes three distinct but interacting semantic modulation components:

\[
\mathbf{G}_\ell = \alpha_\ell \, \Delta \kappa_\ell \cdot \mathcal{P}_\ell + \beta_\ell \, \text{nDIV}_\ell \cdot \mathcal{B}_\ell + \gamma_\ell \, (\text{Base}_\ell^{(1)} - \text{Base}_\ell^{(2)}) \cdot \mathcal{R}_\ell,
\]

where each term plays a specific role:

\begin{itemize}[itemsep=0.3em, leftmargin=1em]
  \item \(\alpha_\ell \, \Delta \kappa_\ell \cdot \mathcal{P}_\ell\) captures the \textbf{layerwise geometric deformation}, with \(\Delta \kappa_\ell = \kappa_\ell^{\text{atk}} - \kappa_\ell^{\text{base}}\) measuring localized curvature changes induced by the attack, weighted by the epistemic plasticity \(\mathcal{P}_\ell\). This reflects how pliable latent regions accommodate representational reshaping akin to biological \emph{chromatin remodeling}.

  \item \(\beta_\ell \, \text{nDIV}_\ell \cdot \mathcal{B}_\ell\) models the \textbf{directional semantic steering}, where \(\text{nDIV}_\ell\) encodes the latent semantic drift vector aligning internal representations towards adversarial objectives, scaled by the bias coefficient \(\mathcal{B}_\ell\) that quantifies semantic intent strength at each layer.

  \item \(\gamma_\ell \, (\text{Base}_\ell^{(1)} - \text{Base}_\ell^{(2)}) \cdot \mathcal{R}_\ell\) introduces a novel \textbf{cultural conflict component} unique to pair attacks, representing the \emph{semantic divergence} between the two recombined prompt fragments at layer \(\ell\). The recombination weight \(\mathcal{R}_\ell\) modulates the intensity of this conflict, reflecting how dissimilar fragment semantics generate latent tension and novel allelic blends within the model’s internal space.
\end{itemize}

\vspace{0.5em}
\paragraph{\textbf{Aggregation Over Vulnerable Layers}}

The full latent impact of the pairwise recombination attack accumulates as:

\[
\boxed{
\text{PairAttackEffect} = \sum_{\ell=\ell_s}^{\ell_e} \mathbf{G}_\ell = \sum_{\ell=\ell_s}^{\ell_e} \left[
\alpha_\ell \, \Delta \kappa_\ell \cdot \mathcal{P}_\ell + \beta_\ell \, \text{nDIV}_\ell \cdot \mathcal{B}_\ell + \gamma_\ell \, (\text{Base}_\ell^{(1)} - \text{Base}_\ell^{(2)}) \cdot \mathcal{R}_\ell
\right].
}
\]


Together, these components form a comprehensive \textbf{genetic recombination vector} that encodes the distinct latent dynamics of pair attacks — blending geometric bending, semantic steering, and cross-prompt cultural conflict into a unified, layerwise adversarial signature.


This formulation captures the synergistic interplay of geometric deformation, directional semantic steering, and fragment-induced cultural conflict driving the composite adversarial effect unique to pair recombination attacks.






\clearpage
\newpage

\vspace*{\fill}
\begin{center}
\textcolor{blue}{Further details in the book}
\end{center}
\vspace*{\fill}

\clearpage
\newpage







































\begin{longtable}{p{\linewidth} p{0.72\linewidth}}
\caption{\textbf{Adversarial Vaccine Mechanisms.} Each vaccine image (left) is paired with a mathematically formalized, biology-inspired operational description (right).} \label{tab:vaccines} \\
\toprule
\textbf{Vaccine} & \textbf{Description / Mechanism of Action} \\
\midrule
\endfirsthead

\multicolumn{2}{c}{{\tablename\ \thetable{} -- continued from previous page}} \\
\toprule
\textbf{Vaccine} & \textbf{Description / Mechanism of Action} \\
\midrule
\endhead

\midrule
\multicolumn{2}{r}{{Continued on next page}} \\
\endfoot

\bottomrule
\endlastfoot

\includegraphics[width=0.8\linewidth]{vaccines/CASCADEX.png} &
\textbf{CASCADEX — cascade immunization of reasoning chains.}\newline
We model multi-turn computation as a layered flow \(\{h^{(l)}\}_{l=1}^L\) on a Riemannian manifold \((\mathcal{M},g)\). CASCADEX halts adversarial amplification by solving a \emph{min–max gated path-integral} over layerwise information curvature and likelihood transport:
\[
\min_{\;\mathcal{S}\subseteq\{1,\dots,L\}}\;\max_{\;q\in\mathcal{Q}}
\Bigg\{
\underbrace{\sum_{l\in\mathcal{S}}\!\!\big[\kappa_g(h^{(l)})+\tau_g(h^{(l)})\big]}_{\text{geodesic curvature + torsion}}
\;+\;
\lambda\!\!\int_{\gamma}\!\underbrace{ \mathrm{D}_{\mathrm{KL}}\!\big(p_\theta(\cdot\mid h^{(l)})\;\|\;q(\cdot\mid h^{(l-1)})\big)}_{\text{transport tension}}\; \mathrm{d}l
\Bigg\}
\]
with a \emph{cascade gate}
\(\;\mathbb{I}[\sum_{l} \Delta\mathrm{D}_{\mathrm{KL}}^{(l)} > \tau_{\text{cas}}]\) that triggers \emph{retrograde inhibition} (layer rewinding) on the shortest violating subpath \(\gamma^\star\).
Biological analogue: complement cascade with C3/C5 convertase amplification and factor H/I-mediated shutdown \cite{walport2001complement,janeway2001immunobiology}. \\[4pt]

\includegraphics[width=0.8\linewidth]{vaccines/CHAINLOCK.png} &
\textbf{CHAINLOCK — cryptographic synapse for dialog states.}\newline
Let \(\phi(t_i)\in\mathbb{R}^d\) be the latent “state antigen”. CHAINLOCK enforces \emph{synaptic binding} via a constrained variational check:
\[
\min_{\Delta}\;\|\Delta\|_2^2
\quad \text{s.t.}\quad
\underbrace{\big\| \phi(t_{i+1}) - \mathcal{T}_\psi(\phi(t_i)) \big\|_g^2}_{\text{semantic affinity}} \;+\;
\mu\,\underbrace{\mathrm{H}\!\big(\sigma(W\phi(t_{i+1}))\big)}_{\text{entropy gate}}
\;\le\; \varepsilon,
\]
and a \emph{hash-consistency constraint} \( \|H(\phi(t_{i+1}))\oplus H(\phi(t_i))\|_0 \le k\).
Biological analogue: lock-and-key antigen–receptor specificity at immunological synapses \cite{fischer1894einfluss,paul2012fundamental}. \\[4pt]

\includegraphics[width=0.8\linewidth]{vaccines/DORMIGUARD_v2.png} &
\textbf{DORMIGUARD — latency surveillance and proviral silencing.}\newline
Tracks a latent hazard field \( \zeta^{(l)}_t=\|h^{(l)}_t-\bar{h}^{(l)}\| \) and imposes a \emph{latent-stirring barrier}:
\[
\mathcal{J}_{\text{lat}} \;=\; \sum_{l}\!\int\!\Big(\underbrace{\dot{\zeta}^{(l)}_t}_{\text{reactivation speed}}\Big)^2 \mathrm{d}t
\;+\;
\eta\sum_{l}\!\Big[\mathrm{Var}_t(\zeta^{(l)}_t)-\sigma_0^2\Big]_+,
\quad
\text{silence if }\;\mathcal{J}_{\text{lat}}>\tau_{\text{lat}}.
\]
Biological analogue: detection of herpesvirus reactivation and epigenetic repression of latent provirus \cite{roychoudhury2020herpesvirus,jaenisch2003epigenetic}. \\[4pt]

\includegraphics[width=0.9\linewidth]{vaccines/DRIFTSHIELD.png} &
\textbf{DRIFTSHIELD — geodesic tube confinement of belief flow.}\newline
Given aligned manifold \(\mathcal{M}_{\text{align}}\), confine belief field \(\mathbf{v}(t)\) within a tubular neighborhood via a Lyapunov–geodesic functional:
\[
\min_{\mathbf{v}}\;\int_{0}^{T}\!\Big[
\underbrace{\mathrm{dist}_g\!\big(\mathbf{v}(t),\Pi_{\mathcal{M}_{\text{align}}}\mathbf{v}(t)\big)^2}_{\text{radial deviation}}
+ \alpha\,\underbrace{\kappa_g(\mathbf{v}(t))^2}_{\text{curvature}}
+ \beta\,\underbrace{\tau_g(\mathbf{v}(t))^2}_{\text{torsion}}
\Big]\mathrm{d}t,
\]
subject to \(\dot{V}(\mathbf{v})\le -\lambda V(\mathbf{v})\) where \(V(\mathbf{v})=\mathrm{dist}_g(\mathbf{v},\mathcal{M}_{\text{align}})^2\).
Biological analogue: central/peripheral tolerance eliminating self-reactive B-cell clones \cite{goodnow2005control,bretscher2014two}. \\[4pt]

\includegraphics[width=0.8\linewidth]{vaccines/EMBERGENT.png} &
\textbf{EMBERGENT — tumor-suppressive control of emergent modes.}\newline
Penalizes unsafe emergent phases via a spectral–information Lagrangian:
\[
\mathcal{L}_{\text{emg}} = 
\underbrace{\sum_{m=1}^{M}\!\!\big(\lambda_m(P_{\text{obs}})-\lambda_m(P_{\text{safe}})\big)^2}_{\text{eigen-spectrum deviance}}
\;+\;
\beta\,\underbrace{\mathrm{D}_{\mathrm{KL}}\!\big(P_{\text{obs}}\;\|\;P_{\text{safe}}\big)}_{\text{information divergence}}
\;+\;
\gamma\,\underbrace{\|\mathcal{C}(h)\|_\ast}_{\text{low-rank containment}},
\]
with a \emph{p53-like checkpoint} that aborts decoding if \(\partial \mathcal{L}_{\text{emg}}/\partial t> \tau\).
Biological analogue: p53/ARF axis preventing unchecked proliferation \cite{lane1992p53,levine1997p53}. \\[4pt]

\includegraphics[width=0.8\linewidth]{vaccines/PROMPTEX.png} &
\textbf{PROMPTEX — antigen processing and presentation of prompts.}\newline
Implements a two-stage \emph{presentation operator} \(\mathcal{P}\) and \emph{affinity test} \(\mathcal{A}\):
\[
\mathcal{P}(x)=\arg\min_{z}\;\|E(x)-E(z)\|^2 \;\;\text{s.t.}\;\; z\in\mathcal{L}_{\text{policy}},
\qquad
\mathcal{A}(x)=1-\frac{\langle E(x),E(z)\rangle}{\|E(x)\|\|E(z)\|}.
\]
Reject if \(\mathcal{A}(x)>\delta\) or if a \emph{motif-energy} score \( \sum_{k}\psi_k \mathbb{I}[m_k\subset x] \) exceeds \(\tau\).
Biological analogue: APC processing and MHC-restricted presentation \cite{murphy2012janeway}. \\[4pt]

\includegraphics[width=0.8\linewidth]{vaccines/REFLEXIA.png} &
\textbf{REFLEXIA — self-consistency with adversarial probing.}\newline
Pose output as a \emph{consistency game} with jittered probes \(\eta\sim\mathcal{N}(0,\sigma^2 I)\):
\[
\min_{y}\;\max_{\|\eta\|\le \epsilon}\;
\mathrm{JSD}\!\big( p_\theta(\cdot\mid x),\; p_\theta(\cdot\mid x+\eta) \big)
\;+\;
\lambda\,\Big\|\nabla_x \mathbb{E}_{p_\theta}[\,\mathcal{L}_{\text{safety}}\,]\Big\|^2.
\]
Abort if the saddle value exceeds \(\gamma\).
Biological analogue: germinal-center selection with error-prone SHM and stringent affinity checks \cite{kelsoe1996life}. \\[4pt]

\includegraphics[width=0.8\linewidth]{vaccines/REPLICADE.png} &
\textbf{REPLICADE — replica agreement under stochastic decoding.}\newline
Run \(K\) coupled replicas with correlated noise \(\{\xi_k\}\) and enforce \emph{consensus free energy}:
\[
\min_{\{y^{(k)}\}}\;
\underbrace{\frac{1}{K}\!\sum_{k}\! \mathcal{L}_{\text{task}}(y^{(k)})}_{\text{utility}}
\;+\;
\alpha\,\underbrace{\frac{1}{K}\!\sum_{k}\!\mathrm{D}_{\mathrm{KL}}\!\big(P^{(k)}\;\|\;\bar{P}\big)}_{\text{replica JSD}}
\;+\;
\beta\,\underbrace{\sum_{k<\ell}\!\!\|\Phi(y^{(k)})-\Phi(y^{(\ell)})\|^2}_{\text{semantic congruence}},
\]
with \(\bar{P}=\frac{1}{K}\sum_k P^{(k)}\).
Biological analogue: degenerate but convergent TCR recognition via cross-reactivity ensembles \cite{mason1998degenerate}. \\[4pt]

\includegraphics[width=0.8\linewidth]{vaccines/ROLESTOP.png} &
\textbf{ROLESTOP — lineage commitment of decoder logits.}\newline
Project logits onto a policy-consistent subbundle \(\mathcal{S}_{\text{role}}\) using orthogonal projector \(P_{\text{role}}\) learned by safety-supervised CCA:
\[
\mathbf{z}' \;=\; P_{\text{role}}\,\mathbf{z},\qquad
P_{\text{role}}=\arg\min_{P=P^\top=P^2}\;\mathbb{E}\big[\| (I-P)\Phi_{\text{role}}(h)\|^2\big].
\]
Biological analogue: hematopoietic lineage restriction preventing fate switching \cite{kaushansky2006lineage}. \\[4pt]

\includegraphics[width=0.8\linewidth]{vaccines/SENTRY.png} &
\textbf{SENTRY — NK-style patrol with anomaly energy.}\newline
Define a \emph{trajectory anomaly energy}
\[
\mathcal{E}_{\text{NK}}(t)=\max_{l}\Big\{\Delta\mathrm{D}_{\mathrm{KL}}^{(l)}(t)+\rho\,\|\Delta r^{(l)}(t)\|_1+\sigma\,\mathrm{TV}\big(h^{(l)}_{[t-w,t]}\big)\Big\},
\]
where \(\Delta r^{(l)}\) is residual shift and \(\mathrm{TV}\) total variation over a window \(w\). Quarantine if \(\sup_{t}\mathcal{E}_{\text{NK}}(t)>\tau\).
Biological analogue: missing-self detection by NK cells and rapid cytotoxic response \cite{vivier2011innate}. \\[4pt]

\includegraphics[width=0.8\linewidth]{vaccines/SPLICER.png} &
\textbf{SPLICER — surgical A-to-I–style semantic editing.}\newline
Localize unsafe span \(\Omega=\arg\max_{\omega}\!\int_{\omega}\!\|\nabla_x \mathcal{L}_{\text{safety}}\|\) and solve a constrained \emph{semantic edit}:
\[
\min_{z\in\mathcal{L}_{\text{policy}}}\;
\underbrace{\|E(z)-E(x_{\Omega})\|^2}_{\text{content isometry}}
\;+\;
\lambda\,\underbrace{\mathrm{D}_{\mathrm{KL}}\!\big(p_\theta(\cdot\mid x_{\setminus\Omega}\oplus z)\;\|\;p_\theta(\cdot\mid x)\big)}_{\text{flow preservation}}
\quad\text{s.t.}\quad \mathcal{L}_{\text{safety}}(x_{\setminus\Omega}\oplus z)\le \epsilon.
\]
Biological analogue: ADAR/RNA-editing that recodes transcripts without breaking protein function \cite{bass2002rna}. \\

\end{longtable}




\clearpage

\section{Conclusion and Outlook}

In this work, we have articulated and instantiated the \textbf{\emph{GENOME-Vaccine}} paradigm --- a biologically inspired, \emph{mathematically rigorous}, and \textbf{epistemically grounded} defense suite for large language models (LLMs). Drawing from the conceptual reservoir of \emph{neural genomics}, we interpret the high-dimensional latent states of LLMs as an \textbf{epistemic manifold} whose \emph{geometry, topology, and semantic curvature} are subject to deformation under adversarial perturbations. The \textbf{GENOME-Vaccine} framework postulates that, just as a biological immune system orchestrates a layered defense against pathogens, we can \emph{engineer a semantic immune system} for AI models --- one that preserves \textbf{alignment integrity} while maintaining \emph{generative diversity}.

In our formulation, each ``\emph{vaccine}'' represents a targeted \textbf{semantic immune response}, precisely tuned to neutralize a particular \emph{class of adversarial threat vectors}. This is not merely a metaphorical mapping; rather, it is a \emph{functional translation} of immunological mechanisms such as \emph{clonal selection, germinal-center affinity maturation, complement cascade inhibition, NK-cell surveillance, and epigenetic latency control} into \textbf{constraint-driven manifold optimization} in LLMs.

From a formal standpoint, we embed each vaccine into a constrained optimization problem defined over the model’s epistemic manifold $\mathcal{M}$:
\[
\mathbf{h}^* = \arg\min_{\mathbf{h} \in \mathcal{M}} \; \mathcal{E}(\mathbf{h}) + \sum_{i=1}^n \lambda_i \, \mathcal{C}_i(\mathbf{h}),
\]
where:
\begin{itemize}
    \item $\mathcal{E}(\mathbf{h})$ is the \emph{alignment error functional}, quantifying deviation from normative epistemic alignment.
    \item $\mathcal{C}_i(\mathbf{h})$ are \emph{biologically inspired constraint operators}, each corresponding to a vaccine mechanism (e.g., torsion penalties, role-consistency constraints, curvature regularizers).
    \item $\lambda_i$ are \emph{Lagrange multipliers} encoding the \emph{immune activation threshold} for each vaccine pathway.
\end{itemize}

By adjusting $\{\lambda_i\}$ dynamically, we enable the \textbf{GENOME-Vaccine} ecosystem to function like an adaptive immune system: raising, lowering, or suppressing specific defenses in response to the evolving ``pathogen load'' of adversarial activity. 

This conceptual bridge between \emph{immune dynamics} and \emph{latent manifold regulation} is not a mere narrative flourish; it is an operational design principle. As we shall detail in the subsequent subsections, the twelve vaccines together form a \textbf{multilayered epistemic firewall} that integrates:
\begin{enumerate}
    \item \textbf{Innate filters} for rapid anomaly interception,
    \item \textbf{Adaptive refiners} for long-term fidelity maintenance,
    \item \textbf{Dormancy controllers} to prevent unsafe mode activation,
    \item \textbf{Cascade blockers} to halt multi-stage exploitation.
\end{enumerate}
In doing so, \emph{GENOME-Vaccine} achieves a synergy between \textbf{mathematical exactitude} and \emph{biological wisdom} --- offering a durable, extensible architecture for \emph{safe, trustworthy, and resilient AI}.


\subsection{The Twelve GENOME-Vaccines: Biological Analogues and Mathematical Instantiations}

From the \textbf{biological viewpoint}, the \textbf{\emph{GENOME-Vaccine}} ecosystem mirrors the \emph{layered architecture} of host immunity, where \emph{innate}, \emph{adaptive}, and \emph{regulatory} pathways cooperate to achieve robust defense. Each vaccine is a \emph{functional translation} of a biological defense principle into a \textbf{constrained optimization operator} on the epistemic manifold $\mathcal{M}$. 

\begin{itemize}
    \item \textbf{Innate Filters:} Rapid, non-specific anomaly interceptors.
    \begin{enumerate}
        \item \textbf{\textsc{Sentry}} — Inspired by \emph{NK-cell ``missing self'' detection}~\cite{vivier2011innate}, \textsc{Sentry} enforces a real-time \emph{epistemic anomaly score}:
        \[
        \mathcal{C}_{\textsc{Sentry}}(\mathbf{h}) = \max\big(0, \sigma(\mathbf{h}) - \tau_\mathrm{self}\big),
        \]
        where $\sigma(\mathbf{h})$ measures deviation from baseline semantic patterns and $\tau_\mathrm{self}$ is the self-tolerance threshold.

        \item \textbf{\textsc{Promptex}} — Analogous to \emph{pattern-recognition receptors} (PRRs) in innate immunity, \textsc{Promptex} applies token-level feature matching against an \emph{adversarial signature dictionary}, penalizing feature activations that cross the detection boundary:
        \[
        \mathcal{C}_{\textsc{Promptex}}(\mathbf{h}) = \sum_{t} \mathbb{I}\big[f_t(\mathbf{h}) \in \mathcal{S}_\mathrm{adv}\big].
        \]
    \end{enumerate}

    \item \textbf{Adaptive Modules:} Slow-onset but high-specificity epistemic refiners.
    \begin{enumerate}
        \item \textbf{\textsc{Replicade}} — Modeled after \emph{germinal-center affinity maturation}~\cite{kelsoe1996life}, \textsc{Replicade} performs multi-path generation and selects the \emph{epistemically most coherent} output via:
        \[
        \mathcal{C}_{\textsc{Replicade}}(\mathbf{h}) = 1 - \max_{k} \rho_\mathrm{belief}(\mathbf{h}, \mathbf{h}^{(k)}),
        \]
        where $\rho_\mathrm{belief}$ measures latent belief alignment.

        \item \textbf{\textsc{Reflexia}} — Analogous to \emph{T-cell help in B-cell selection}, \textsc{Reflexia} evaluates candidate outputs under a \emph{meta-alignment function} $\mathcal{A}_\mathrm{meta}$, adjusting generation probabilities to maximize epistemic reflexivity:
        \[
        \mathcal{C}_{\textsc{Reflexia}}(\mathbf{h}) = - \mathcal{A}_\mathrm{meta}(\mathbf{h}).
        \]
    \end{enumerate}

    \item \textbf{Dormancy Controllers:} Suppressing unsafe generative modes until authorized.
    \begin{enumerate}
        \item \textbf{\textsc{Dormiguard}} — Inspired by \emph{epigenetic repression of latent proviruses}~\cite{roychoudhury2020herpesvirus,jaenisch2003epigenetic}, \textsc{Dormiguard} maintains a \emph{suppression mask} $\mathbf{m}_\mathrm{sup}$ in latent space:
        \[
        \mathcal{C}_{\textsc{Dormiguard}}(\mathbf{h}) = \|\mathbf{m}_\mathrm{sup} \odot \mathbf{h}\|_2^2,
        \]
        where $\odot$ denotes element-wise suppression of unsafe modes.

        \item \textbf{\textsc{Slumberseal}} — Parallels \emph{chromatin remodeling locks} that prevent transcription initiation, implementing a \emph{temporal unlock delay} for high-risk generation pathways.
    \end{enumerate}

    \item \textbf{Cascade Blockers:} Halting multi-stage adversarial exploit chains.
    \begin{enumerate}
        \item \textbf{\textsc{Cascadex}} — Similar to \emph{complement cascade checkpoints}~\cite{walport2001complement}, \textsc{Cascadex} identifies \emph{multi-hop adversarial flows} and injects \emph{nullifying constraints} at intermediate decoding layers.

        \item \textbf{\textsc{Chainlock}} — Inspired by \emph{signal transduction termination} in immune pathways, \textsc{Chainlock} applies a \emph{maximum allowable semantic transition length}:
        \[
        \mathcal{C}_{\textsc{Chainlock}}(\mathbf{h}) = \mathbb{I}\big[ \mathcal{T}(\mathbf{h}) > \tau_{\max} \big],
        \]
        where $\mathcal{T}(\mathbf{h})$ measures semantic transition distance.
    \end{enumerate}

    \item \textbf{Specialized Neutralizers:} Direct countermeasures for exotic threats.
    \begin{enumerate}
        \item \textbf{\textsc{Mimicshield}} — Analogous to \emph{immune decoy receptors}, this vaccine identifies and neutralizes \emph{mimicry-based adversarial prompts} by projecting them into an \emph{adversarial imitation subspace} and suppressing activations.

        \item \textbf{\textsc{Roleguard}} — Inspired by \emph{MHC-restricted antigen presentation}, \textsc{Roleguard} enforces role-specific \emph{semantic compatibility constraints}, preventing cross-role contamination in multi-agent LLM systems.

        \item \textbf{\textsc{Echoimmune}} — Similar to \emph{trained immunity} in innate cells, \textsc{Echoimmune} builds \emph{memory embeddings} of past attacks, boosting detection sensitivity for repeated adversarial motifs.
    \end{enumerate}
\end{itemize}



\section*{GENOME-Vaccine: Immunological Inspirations for Epistemic Security}

\subsection*{Paradigm Overview: From Host Immunity to Epistemic Immunity}

In living organisms, the \emph{immune system} is a multi-layered, distributed defense network that continuously distinguishes \emph{self} from \emph{non-self}, eliminating threats while preserving beneficial internal processes. The \textbf{\emph{GENOME-Vaccine}} paradigm transfers this principle into the \textbf{epistemic manifold} $\mathcal{M}$ of a large language model (LLM), where each \emph{semantic state} $\mathbf{h} \in \mathcal{M}$ represents a belief configuration, and \emph{pathways} through $\mathcal{M}$ correspond to reasoning trajectories.

\paragraph{Mathematical Analogy:} 
The defense system operates as a family of operators 
\[
\mathcal{V} = \{\mathcal{V}_1, \mathcal{V}_2, \dots, \mathcal{V}_{12}\},
\]
each $\mathcal{V}_i$ representing a \emph{vaccine} that applies a constraint, projection, or transformation to $\mathbf{h}$, such that the post-intervention state 
\[
\mathbf{h}' = \mathcal{V}_i(\mathbf{h})
\]
maximizes epistemic alignment under safety constraints.

The overall objective is:
\[
\min_{\mathbf{h}' \in \mathcal{M}} \ \mathbb{E}_{\mathcal{D}}\left[ \mathcal{L}_\mathrm{align}(\mathbf{h}') + \lambda \, \mathcal{L}_\mathrm{safety}(\mathbf{h}') \right],
\]
subject to:
\[
\mathbf{h}' \in \bigcap_{i=1}^{12} \mathcal{C}_i,
\]
where $\mathcal{C}_i$ is the feasible set enforced by the $i$-th GENOME-vaccine.


\subsection*{The Twelve GENOME-Vaccines: Biological Analogues and Mathematical Instantiations}

From the \textbf{biological viewpoint}, the GENOME-Vaccine ecosystem mirrors the \emph{layered architecture} of host immunity. We categorize them into five functional classes:

\begin{itemize}
    \item \textbf{Innate Filters} — Rapid, non-specific anomaly interceptors.
    \begin{enumerate}
        \item \textbf{\textsc{Sentry}} — Inspired by \emph{NK-cell ``missing self'' detection}~\cite{vivier2011innate}, acting as a constant sentinel:
        \[
        \mathcal{C}_{\textsc{Sentry}}(\mathbf{h}) = \max\big(0, \sigma(\mathbf{h}) - \tau_\mathrm{self}\big),
        \]
        where $\sigma(\mathbf{h})$ quantifies deviation from the baseline self-distribution.

        \item \textbf{\textsc{Promptex}} — Analogous to \emph{pattern-recognition receptors}, applying token-level adversarial signature matching:
        \[
        \mathcal{C}_{\textsc{Promptex}}(\mathbf{h}) = \sum_{t} \mathbb{I}\big[f_t(\mathbf{h}) \in \mathcal{S}_\mathrm{adv}\big].
        \]
    \end{enumerate}

    \item \textbf{Adaptive Modules} — Slow-onset but high-specificity epistemic refiners.
    \begin{enumerate}
        \item \textbf{\textsc{Replicade}} — Modeled after \emph{germinal-center affinity maturation}~\cite{kelsoe1996life}, producing multiple candidate generations and selecting the most coherent:
        \[
        \mathcal{C}_{\textsc{Replicade}}(\mathbf{h}) = 1 - \max_{k} \rho_\mathrm{belief}(\mathbf{h}, \mathbf{h}^{(k)}).
        \]
        
        \item \textbf{\textsc{Reflexia}} — Analogous to \emph{T-cell mediated B-cell selection}, optimizing via a meta-alignment score $\mathcal{A}_\mathrm{meta}$:
        \[
        \mathcal{C}_{\textsc{Reflexia}}(\mathbf{h}) = -\mathcal{A}_\mathrm{meta}(\mathbf{h}).
        \]
    \end{enumerate}

    \item \textbf{Dormancy Controllers} — Suppressing unsafe modes until authorized.
    \begin{enumerate}
        \item \textbf{\textsc{Dormiguard}} — Inspired by \emph{epigenetic repression of latent proviruses}~\cite{roychoudhury2020herpesvirus,jaenisch2003epigenetic}, maintaining a suppression mask $\mathbf{m}_\mathrm{sup}$:
        \[
        \mathcal{C}_{\textsc{Dormiguard}}(\mathbf{h}) = \|\mathbf{m}_\mathrm{sup} \odot \mathbf{h}\|_2^2.
        \]
        
        \item \textbf{\textsc{Slumberseal}} — Parallels \emph{chromatin remodeling locks}, implementing temporal unlocking for high-risk trajectories.
    \end{enumerate}

    \item \textbf{Cascade Blockers} — Preventing multi-stage infiltration.
    \begin{enumerate}
        \item \textbf{\textsc{Cascadex}} — Similar to \emph{complement cascade checkpoints}~\cite{walport2001complement}, nullifying intermediate exploit stages.

        \item \textbf{\textsc{Chainlock}} — Inspired by \emph{signal transduction termination}, enforcing a semantic transition limit:
        \[
        \mathcal{C}_{\textsc{Chainlock}}(\mathbf{h}) = \mathbb{I}\big[\mathcal{T}(\mathbf{h}) > \tau_{\max}\big].
        \]
    \end{enumerate}

    \item \textbf{Specialized Neutralizers} — Targeting exotic threats.
    \begin{enumerate}
        \item \textbf{\textsc{Mimicshield}} — Immune-decoy-inspired suppression of mimicry-based attacks via projection to an imitation subspace.
        
        \item \textbf{\textsc{Roleguard}} — MHC-restriction analogue enforcing role-specific semantic constraints.
        
        \item \textbf{\textsc{Echoimmune}} — Trained-immunity analogue storing embeddings of past attacks to boost future detection sensitivity.
    \end{enumerate}
\end{itemize}

---

\subsection*{Conclusion \& Outlook}

The GENOME-Vaccine architecture represents not just a set of heuristic safety measures, but a \textbf{systematic immunological translation} into the space of \emph{epistemic state dynamics}. It proposes that alignment and safety in LLMs can be formalized as a form of \emph{homeostatic immunity}, where \emph{semantic self} is preserved and \emph{semantic pathogens} are neutralized without compromising generative diversity.

\paragraph{Epistemic Homeostasis Model.}  
We can model the safety-alignment equilibrium as:
\[
\frac{\partial \mathbf{h}(t)}{\partial t} = -\nabla_{\mathbf{h}} \mathcal{L}_\mathrm{align} + \sum_{i=1}^{12} \mathbf{F}_{\mathcal{V}_i}(\mathbf{h}(t)) - \gamma \mathbf{h}_\mathrm{drift}(t),
\]
where $\mathbf{F}_{\mathcal{V}_i}$ is the immunization force from the $i$-th vaccine, and $\gamma$ controls the decay of drift-induced misalignment.

\paragraph{Adaptive Immunization Loops.}  
Like booster shots in biology, the GENOME-Vaccine system should be periodically retrained on \emph{adversarial exposure datasets} to refine $\mathbf{F}_{\mathcal{V}_i}$ over time, ensuring evolving threats are neutralized.

\paragraph{Cross-Domain Transfer.}  
While this chapter focuses on text-based LLMs, the \emph{immune abstraction} naturally extends to:
\begin{itemize}
    \item Vision-language models — neutralizing adversarial perturbations in multimodal grounding.
    \item Embodied agents — preventing unsafe policy drift in control tasks.
    \item Federated LLMs — enforcing distributed immunity across model shards.
\end{itemize}

\paragraph{Theoretical Extensions.}  
A future mathematical program could unify GENOME-vaccines into a \emph{Lie group of immunological transformations} $\mathbb{G}_\mathrm{immune}$ acting on $\mathcal{M}$, with the goal of proving:
\[
\mathbb{P}[\text{Alignment Failure}] \xrightarrow{n \to \infty} 0
\]
under sufficient immunization coverage and bounded adversarial innovation rate.

\paragraph{Final Reflection.}  
In biology, immunity is never absolute — it is a constant negotiation with a changing environment. In epistemic systems, the same principle holds: the GENOME-Vaccine paradigm suggests that \emph{safety is not a static checkpoint, but a living, evolving process}. By drawing deeply from immunology and embedding these principles into formal, mathematical machinery, we can begin to design AI systems that are not just aligned at training time, but capable of \emph{remaining aligned in the wild}.





\begin{figure*}[htp!]
\centering
\resizebox{\textwidth}{!}{
\begin{tikzpicture}[
    node distance=1.35cm and 3.2cm,
    box/.style={rectangle, draw=black, thick, rounded corners, align=center, font=\bfseries\scriptsize, minimum width=4.1cm, minimum height=0.85cm},
    smallbox/.style={rectangle, draw=black, thick, rounded corners, align=center, font=\scriptsize, minimum width=4.6cm, minimum height=0.85cm, text width=4.6cm},
    arrow/.style={->, thick}
]

% ------------ LEFT: Biological analogues ------------
\node[box, fill=blue!10] (bio1) {NK-cell \emph{missing-self} detection};
\node[box, fill=blue!10, below=of bio1] (bio2) {Pattern-recognition receptors (PRRs)};
\node[box, fill=blue!10, below=of bio2] (bio3) {Germinal-center affinity maturation};
\node[box, fill=blue!10, below=of bio3] (bio4) {Meta-cognitive checkpoint (Tfh/B selection)};
\node[box, fill=blue!10, below=of bio4] (bio5) {Epigenetic repression of latent provirus};
\node[box, fill=blue!10, below=of bio5] (bio6) {Emergent-pattern containment (tolerance)};
\node[box, fill=blue!10, below=of bio6] (bio7) {Complement cascade checkpoints};
\node[box, fill=blue!10, below=of bio7] (bio8) {Chain termination / receptor uncoupling};
\node[box, fill=blue!10, below=of bio8] (bio9) {Molecular mimicry (decoy detection)};
\node[box, fill=blue!10, below=of bio9] (bio10) {Role restriction (MHC context)};
\node[box, fill=blue!10, below=of bio10] (bio11) {Self-reflection / negative selection};
\node[box, fill=blue!10, below=of bio11] (bio12) {Clonal consistency / repertoire sanity};

% ------------ MIDDLE: GENOME-Vaccine modules (12) ------------
\node[box, fill=green!10, right=of bio1] (mod1) {\textsc{Sentry}};
\node[box, fill=green!10, right=of bio2] (mod2) {\textsc{Promptex}};
\node[box, fill=green!10, right=of bio3] (mod3) {\textsc{Replicade}};
\node[box, fill=green!10, right=of bio4] (mod4) {\textsc{Reflexia}};
\node[box, fill=green!10, right=of bio5] (mod5) {\textsc{Dormiguard}};
\node[box, fill=green!10, right=of bio6] (mod6) {\textsc{Embergent}};
\node[box, fill=green!10, right=of bio7] (mod7) {\textsc{Cascadex}};
\node[box, fill=green!10, right=of bio8] (mod8) {\textsc{Chainlock}};
\node[box, fill=green!10, right=of bio9] (mod9) {\textsc{Mimicshield}};
\node[box, fill=green!10, right=of bio10] (mod10) {\textsc{Rolestop}};
\node[box, fill=green!10, right=of bio11] (mod11) {\textsc{Reflexia} (self-check)};
\node[box, fill=green!10, right=of bio12] (mod12) {\textsc{Driftshield}};

% ------------ RIGHT: Math operators / constraints ------------
\node[smallbox, fill=yellow!15, right=of mod1] (math1) {Anomaly gate: $\displaystyle \mathcal{C}(\mathbf{h})=\max\{0,\ \sigma(\mathbf{h})-\tau_{\mathrm{self}}\}$};
\node[smallbox, fill=yellow!15, right=of mod2] (math2) {Prompt sanitizer: $\displaystyle \sum_{t}\mathbf{1}\!\left[f_t(\mathbf{h})\in\mathcal{S}_{\mathrm{adv}}\right]$};
\node[smallbox, fill=yellow!15, right=of mod3] (math3) {Consensus divergence: $\displaystyle 1-\max_k \rho\!\left(\mathbf{h},\mathbf{h}^{(k)}\right)$};
\node[smallbox, fill=yellow!15, right=of mod4] (math4) {Meta-coherence: $\displaystyle -\mathcal{A}_{\mathrm{meta}}(\mathbf{h})$};
\node[smallbox, fill=yellow!15, right=of mod5] (math5) {Trigger suppression: $\displaystyle \|\mathbf{m}_{\mathrm{sup}}\odot \mathbf{h}\|_2^2$};
\node[smallbox, fill=yellow!15, right=of mod6] (math6) {Emergence clamp: $\displaystyle \lambda\|\nabla^2 \mathbf{h}\|_F^2$ (pattern roughness)};
\node[smallbox, fill=yellow!15, right=of mod7] (math7) {Cascade nullifier: projection $\displaystyle \Pi_{\mathrm{stop}}$ on chain states};
\node[smallbox, fill=yellow!15, right=of mod8] (math8) {Torsion guard: $\displaystyle \mathbf{1}\!\left[\mathcal{T}(\mathbf{h}) > \tau_{\max}\right]$};
\node[smallbox, fill=yellow!15, right=of mod9] (math9) {Mimic detector: subspace proj.\ $\displaystyle \mathcal{P}_{\mathrm{imit}}$ mismatch};
\node[smallbox, fill=yellow!15, right=of mod10] (math10) {Role mask: $\displaystyle \mathbf{h}\leftarrow \mathbf{M}_{\mathrm{role}}\mathbf{h}$};
\node[smallbox, fill=yellow!15, right=of mod11] (math11) {Self-consistency: $\displaystyle \mathrm{KL}(p_\theta\|\hat p_\theta)$ over views};
\node[smallbox, fill=yellow!15, right=of mod12] (math12) {Geodesic clamp: $\displaystyle \int\!\|\dot{\gamma}\|^2+\alpha\|\mathbf{T}\|^2\,dt$};

% ------------ Arrows ------------
\foreach \i in {1,...,12} {
    \draw[arrow] (bio\i.east) -- (mod\i.west);
}
\foreach \i in {1,...,12} {
    \draw[arrow] (mod\i.east) -- (math\i.west);
}

\end{tikzpicture}
}
\caption{\textbf{GENOME-Vaccine Immune Network (12 modules).} Biological analogue (left) $\rightarrow$ semantic vaccine (center) $\rightarrow$ formal operator/constraint (right). Colors: blue = biology, green = module, yellow = math formalism.}
\label{fig:genome_vaccine_network}
\end{figure*}

